\section{Conclusion}

This paper presented City1 v2 as a reproducible open-data baseline for officially calibrated intra-urban proxy population surface modeling under missing cell-level truth in Kazakhstan. The system combines OpenStreetMap-derived urban features, weak supervision, supervised learning, and calibration to official city population totals within a frozen and fully traceable workflow. Rather than claiming exact reconstruction of census counts at the grid-cell level, the paper positions City1 v2 as a bounded analytical baseline for settings in which fine-resolution ground-truth population labels are unavailable.

The frozen evidence stack supports three main conclusions. First, official-total calibration is not a peripheral adjustment but a core component of a coherent city-level proxy population surface. Second, the selected production identity of the system---\emph{random forest} on a \emph{500 m} grid in \emph{calibrated-only} mode---is supported by the validation results reported in the manuscript. Third, no single evaluation layer is sufficient on its own in this problem setting; the contribution of City1 v2 therefore lies not only in the surface itself, but also in the multi-level validation design spanning transferability, within-city robustness, partial internal administrative comparison, external structural benchmarking, ablation, and qualitative plausibility assessment. Supplementary Table~S1 makes that evidence logic explicit by showing what each layer supports and what it cannot be asked to prove.

At the same time, the paper maintains deliberately bounded claims. City1 v2 does not provide true grid-level census observations, and its output should not be interpreted as ground-truth population per cell. The system remains dependent on weak supervision, variable OSM quality, and official-total calibration, while district-level comparison and external benchmark agreement provide support that is useful but necessarily partial. These constraints are not incidental limitations hidden behind the results; they define the practical and methodological setting in which the baseline is intended to be used.

Within those bounds, City1 v2 offers a useful contribution for data-scarce urban analysis. It provides a transparent, reproducible, and empirically supported reference point for intra-urban spatial work in contexts where detailed census microdata are unavailable but open geospatial data and official city totals exist. More broadly, the manuscript contributes a disciplined paper package in which figures, tables, and narrative claims are tied to a frozen evidence stack rather than to moving experimental targets.

Future extensions, including uncertainty modeling, additional remote-sensing or accessibility signals, broader multimodal feature integration, or expanded comparative benchmarking, belong to a subsequent version of the system rather than to the present paper. The contribution of this manuscript is therefore intentionally narrow but solid: City1 v2 establishes a reproducible calibrated proxy population surface baseline with multi-level validation for intra-urban analysis in data-scarce settings.
